% 4_medio.tex

% Investigue sobre el funcionamiento del medio inalámbrico seleccionado y sus limitaciones para transmisión de información digital.

\section{Funcionamiento y limitaciones del medio}

En el espacio libre, los campos eléctricos y magnéticos no los encierra ninguna estructura de confinamiento, por lo que pueden tener cualquier magnitud y dirección, las cuales se determinan por el dispositivo que las genere, por ejemplo, una antena \cite{Hyatt}.
El medio inalámbrico seleccionado para transmitir la información digital es la radiofrecuencia, la cual consiste en la propagación de ondas por medio del espacio libre \cite{Hyatt}, donde las ondas transportan la información de los archivos de audio.
Estas son emitidas y captadas por las antenas incorporadas en el emisor y el receptor.
La función que cumplen las antenas es irradiar potencia en todas direcciones; a estas se les conoce como antenas isotrópicas, de forma que la señal sea captada por una antena receptora.
Las ondas viajan a la velocidad de la luz, donde su potencia depende de la potencia de transmisión, y al existir pérdidas el receptor recibe las ondas electromagnéticas con una potencia menor.

\subsection{Limitaciones del medio de transmisión}

Dado que la comunicación se da en un medio no guiado, la información digital está expuesta a sufrir los fenómenos típicos del aire: atenuación, obstáculos que bloquean la señal, ruido electromagnético, reflexión de la onda sobre otros objetos y dispersión \cite{wireless}.
A continuación, se explican algunos de ellos.

En primer lugar, la atenuación de la señal describe cómo disminuye la potencia recibida con la distancia entre el transmisor y el receptor.
Se debe a la disipación de la potencia radiada por el transmisor, así como a los efectos del canal de propagación.

En segundo lugar, una señal transmitida por medio de un canal inalámbrico suele experimentar variaciones aleatorias debido a la obstrucción causada por objetos en su trayectoria, lo que da lugar a variaciones aleatorias en la potencia recibida a una distancia determinada.
Dichas variaciones también se deben a cambios en las superficies reflectantes y a la dispersión de la señal por objetos.
El \textit{shadowing}, conocido como el efecto sombra, es la atenuación causada por obstáculos entre el transmisor y el receptor que absorben la señal transmitida.
Cuando la atenuación es significativa, el objeto puede bloquear la señal directamente.
Esta afectación en la señal no depende directamente de la distancia entre el transmisor y el receptor, sino de la cantidad y tamaño de los objetos que se encuentren en línea de vista.
La variación debida a este efecto se produce a distancias proporcionales a la longitud del objeto que la obstruye (\SIrange{10}{100}{\meter} en exteriores, y menos en interiores) \cite{wireless}.

Luego, se tienen la reflexión, la difracción y la dispersión.
Estos fenómenos suceden cuando elementos del entorno interactúan con los elementos del transmisor o receptor.
Los componentes de la señal que surgen debido a estos objetos se denominan componentes de trayectos múltiples.
Los diferentes componentes de trayectos múltiples llegan al receptor con distintos retardos y desfases \cite{wireless}.

Una vez conocido el funcionamiento del canal y los fenómenos que genera, se pueden enumerar las limitaciones que presenta el método de transmisión de la información digital: ancho de banda, capacidad del canal, potencia de transmisión, distancia emisor-receptor y ruido electromagnético.

\subsection{Uso de bandas no licenciadas}

Según \cite{SUTEL}, el espectro radioeléctrico es un bien natural de dominio público propio de la nación costarricense, por lo que no puede salir definitivamente del dominio del Estado.
Por esta razón, se permite la comunicación por medio de radiofrecuencia únicamente en ciertas bandas no licenciadas, como la banda de \SI{2.4}{\giga\hertz} y la de \SI{900}{\mega\hertz}, entre otras.
Se seleccionó la banda de \SI{2.4}{\giga\hertz} porque corresponde a la frecuencia de operación del \textit{transceiver} disponible en el laboratorio.
Además, de acuerdo con el modelo de atenuación de Friis visto en el curso, la atenuación aumenta con la frecuencia de la banda de transmisión:
\begin{equation}
    L = \left(\frac{4 \pi d f}{c}\right)^2,
\end{equation}
lo cual implica que se tiene mayor atenuación respecto a otras opciones, pero está dentro de un rango aceptable para las distancias propuestas sobre las que se realiza la comunicación punto a punto.

El \textit{transceiver} NRF24L01, disponible en el laboratorio, opera en la banda ISM de \SI{2.4}{\giga\hertz}, la cual es de uso libre y no requiere licencias para su utilización.
Sin embargo, esta banda es compartida con tecnologías como WiFi y Bluetooth, por lo que pueden presentarse interferencias que afecten la calidad de la comunicación, provoquen pérdidas de paquetes y reduzcan el alcance efectivo del enlace inalámbrico.

